\documentclass[11pt]{article}
\usepackage[margin=1in]{geometry}
\usepackage{amsmath,amssymb,amsthm,enumitem,hyperref}
\hypersetup{colorlinks=true,linkcolor=blue,urlcolor=blue,citecolor=blue}
\newtheorem{theorem}{Theorem}
\newcommand{\HH}{\mathcal H}
\title{A negative leading coefficient in the third LIS iterate\\\large MathDB 317479}
\author{Principia Math In Progress}
\date{September 18, 2026}
\begin{document}
\maketitle
\begin{abstract}
The longest-increasing-subsequence polynomials are not infinitely
$q$-log-convex. Their third log-convexity iterate at index $5$
has degree $34$ and leading coefficient $-32$. Only the top
three coefficients of $P_2,\ldots,P_8$ are needed, giving a short
exact disproof of Chen's Conjecture 1.4.
\end{abstract}

\section*{Task 1. Full disproof}
Let $P_{n,k}$ count permutations of $\{1,\ldots,n\}$ whose
longest increasing subsequence has length $k$, and put
\[
 P_n(x)=\sum_{k=1}^nP_{n,k}x^k.
\]
Following \cite[p.~5]{Chen}, use the centered operator
\[
 (\HH A)_n=A_{n-1}A_{n+1}-A_n^2.
\]
Infinite $q$-log-convexity means coefficientwise nonnegativity
of every iterate. The letter $x$ plays the role of $q$.

\begin{theorem}
The third iterate satisfies
\[
 (\HH^3P)_5(x)=-32x^{34}+\text{a polynomial of degree at most }33.
                                                               \tag{1}
\]
Consequently Chen's infinite $q$-log-convexity conjecture is false.
\end{theorem}

\begin{proof}
Write the top three coefficients as
\[
 P_n(x)=x^n+a_nx^{n-1}+b_nx^{n-2}
        +\text{lower-degree terms}.
\]
The needed values, also printed in \cite[pp.~2--3]{Chen}, are
\[
\begin{array}{c|rrrrrrr}
n&2&3&4&5&6&7&8\\\hline
a_n&1&4&9&16&25&36&49\\
b_n&0&1&13&61&181&421&841
\end{array}                                                    \tag{2}
\]
These counts can be obtained without enumerating all permutations.
The Robinson--Schensted correspondence gives
$P_{n,k}=\sum_{\lambda\vdash n,\,\lambda_1=k}(f^\lambda)^2$,
where $f^\lambda$ counts standard Young tableaux of shape $\lambda$.
For $k=n-1$ there is just the shape $(n-1,1)$, so
$a_n=(n-1)^2$. For $n\ge4$ and $k=n-2$ the shapes are
$(n-2,2)$ and $(n-2,1,1)$. The hook-length formula gives
\[
 b_n=\left(\frac{n(n-3)}2\right)^2
       +\left(\frac{(n-1)(n-2)}2\right)^2.
\]
The two smaller values are $b_2=0$ and $b_3=1$. This proves (2).

Put $B_n=(\HH P)_n$ and $C_n=(\HH^2P)_n$.
The leading terms of degree $2n$ cancel in $B_n$. The next
coefficient is
$a_{n-1}+a_{n+1}-2a_n=2$. Thus, for $3\le n\le7$,
\[
 B_n(x)=2x^{2n-1}+v_nx^{2n-2}+\text{lower-degree terms},
\]
where multiplication gives
\[
 v_n=b_{n-1}+b_{n+1}+a_{n-1}a_{n+1}-a_n^2-2b_n.
\]
The complete required table is
\[
 \begin{array}{c|rrrrr}
 n&3&4&5&6&7\\\hline
 v_n&4&19&41&71&109
 \end{array}.                                                 \tag{3}
\]
In $C_n=B_{n-1}B_{n+1}-B_n^2$, the terms of degree $4n-2$
again cancel. The coefficient of degree $4n-3$ is
\[
 c_n=2(v_{n-1}+v_{n+1}-2v_n).
\]
Using (3),
\[
 C_4=14x^{13}+\cdots,\qquad
 C_5=16x^{17}+\cdots,\qquad
 C_6=16x^{21}+\cdots.                                        \tag{4}
\]
Therefore $C_4C_6-C_5^2$ has degree $34$ and leading
coefficient $14\cdot16-16^2=-32$, proving (1).
\end{proof}

The witness uses only $P_2,\ldots,P_8$; it involves no boundary
convention for $P_0$. It also gives a failure at an ordinary
positive parameter. Direct evaluation of the same polynomials yields
\[
 (\HH^3P)_5(6)=-8449005348498088692649082880000<0.              \tag{5}
\]
The exact numeric iterate table is included in the verification output.

\section*{Verification and scope}
The verifier enumerates all $46232$ permutations of sizes $2$ through
$8$ and checks each longest-increasing-subsequence length by two
algorithms. It independently reconstructs all seven polynomials by
the hook-length formula, checks the full polynomial iterates by two
exact multiplication methods, and verifies (5). No floating point
is used. Review is by the authoring agent only.

The source's reported computational evidence for third-order
$q$-log-convexity on p.~5 is incompatible with (1).
The definition and printed input polynomials themselves agree with
the independent enumeration. This is a counterexample to the intended
conjecture, not a result based on a missing endpoint or a transcription
error. It does not decide ordinary $q$-log-convexity, second-order
$q$-log-convexity, or infinite log-concavity of the coefficient
sequence of each individual $P_n$.

\section*{Tasks 2--5. Follow-up directions}
\begin{enumerate}[start=2,leftmargin=*]
\item Decide whether $(P_n(x))_{n\ge1}$ is second-order
$q$-log-convex.
\item Determine the positive parameters $x$ for which the numerical
sequence $(P_n(x))_{n\ge1}$ is infinitely log-convex.
\item Locate all negative coefficients in higher iterates and
determine whether each fixed iterate is eventually coefficientwise
nonnegative as $n$ increases.
\item Derive formulas for the leading nonzero coefficients of
$\HH^rP$ for general $r$, using tableaux close to a single row.
\end{enumerate}

\begin{thebibliography}{9}
\raggedright
\bibitem{Chen} W. Y. C. Chen,
\emph{Log-concavity and $q$-Log-convexity Conjectures on the
Longest Increasing Subsequences of Permutations},
\href{https://arxiv.org/abs/0806.3392v2}{arXiv:0806.3392v2},
June 22, 2008. Conjecture 1.4, pp.~1--5.
\bibitem{BLS} M. B\'ona, M.-L. Lackner and B. E. Sagan,
\emph{Longest Increasing Subsequences and Log Concavity},
Ann. Comb. 21 (2017), 535--549,
\href{https://doi.org/10.1007/s00026-017-0365-x}{doi:10.1007/s00026-017-0365-x}.
Sections 6.3--6.4 distinguish related conjectures.
\bibitem{DB} MathDB, \href{https://mathdb.com/p/317479}{Problem 317479},
accessed September 18, 2026.
\end{thebibliography}
\end{document}
